% ========================================================
% BÀI TẬP TOÁN 5 - KẾT NỐI TRI THỨC VỚI CUỘC SỐNG
% Bài 5: Ôn tập các phép tính với phân số
% Định dạng: BTTeX Part III - Trắc nghiệm điền số ngắn
% ========================================================

% Câu 1
\begin{ex}
  Tính kết quả: $\frac{3}{5} \times \frac{5}{7}$. Phân số tối giản có tử số là bao nhiêu?
  \shortans{3}
  \loigiai{
    $\frac{3 \times 5}{5 \times 7} = \frac{3}{7}$. Tử số là 3.
  }
\end{ex}

% Câu 2
\begin{ex}
  Tính kết quả: $\frac{4}{9} \times \frac{3}{8}$. Phân số tối giản có mẫu số là bao nhiêu?
  \shortans{6}
  \loigiai{
    $\frac{4 \times 3}{9 \times 8} = \frac{12}{72} = \frac{1}{6}$. Mẫu số là 6.
  }
\end{ex}

% Câu 3
\begin{ex}
  Tính giá trị của phép chia: $\frac{5}{8} : \frac{5}{16}$. Kết quả là bao nhiêu?
  \shortans{2}
  \loigiai{
    $\frac{5}{8} : \frac{5}{16} = \frac{5}{8} \times \frac{16}{5} = 2$.
  }
\end{ex}

% Câu 4
\begin{ex}
  Tính giá trị của phép nhân: $\frac{7}{12} \times 36$. Kết quả là bao nhiêu?
  \shortans{21}
  \loigiai{
    $\frac{7 \times 36}{12} = 7 \times 3 = 21$.
  }
\end{ex}

% Câu 5
\begin{ex}
  Tính giá trị phép chia: $45 : \frac{5}{3}$. Kết quả là bao nhiêu?
  \shortans{27}
  \loigiai{
    $45 : \frac{5}{3} = 45 \times \frac{3}{5} = 27$.
  }
\end{ex}

% Câu 6
\begin{ex}
  Tính nhanh: $\frac{3}{7} \times \frac{4}{9} + \frac{3}{7} \times \frac{5}{9}$. Tử số của kết quả tối giản là bao nhiêu?
  \shortans{3}
  \loigiai{
    $\frac{3}{7} \times (\frac{4}{9} + \frac{5}{9}) = \frac{3}{7} \times 1 = \frac{3}{7}$. Tử số là 3.
  }
\end{ex}

% Câu 7
\begin{ex}
  Tính giá trị: $\frac{2}{3} \times \frac{3}{4} \times \frac{4}{5} \times \frac{5}{6}$. Phân số tối giản có mẫu số là bao nhiêu?
  \shortans{3}
  \loigiai{
    Rút gọn các thừa số liên tiếp: $\frac{2}{6} = \frac{1}{3}$. Mẫu số là 3.
  }
\end{ex}

% Câu 8
\begin{ex}
  Tìm $x$ biết: $x \times \frac{3}{4} = \frac{9}{8}$. Mẫu số của phân số tối giản $x$ là bao nhiêu?
  \shortans{2}
  \loigiai{
    $x = \frac{9}{8} : \frac{3}{4} = \frac{9}{8} \times \frac{4}{3} = \frac{3}{2}$. Mẫu số là 2.
  }
\end{ex}

% Câu 9
\begin{ex}
  Tìm $x$ biết: $x : \frac{2}{5} = 15$. Kết quả $x$ là bao nhiêu?
  \shortans{6}
  \loigiai{
    $x = 15 \times \frac{2}{5} = 6$.
  }
\end{ex}

% Câu 10
\begin{ex}
  Tìm $x$ biết: $\frac{4}{7} : x = \frac{2}{21}$. Kết quả $x$ là bao nhiêu?
  \shortans{6}
  \loigiai{
    $x = \frac{4}{7} : \frac{2}{21} = \frac{4}{7} \times \frac{21}{2} = 6$.
  }
\end{ex}

% Câu 11
\begin{ex}
  Một hình chữ nhật có chiều dài $\frac{4}{5}\text{ m}$, chiều rộng $\frac{3}{8}\text{ m}$. Diện tích hình chữ nhật đó có tử số khi rút gọn về phân số tối giản là bao nhiêu?
  \shortans{3}
  \loigiai{
    Diện tích là $\frac{4}{5} \times \frac{3}{8} = \frac{12}{40} = \frac{3}{10}\text{ m}^2$. Tử số là 3.
  }
\end{ex}

% Câu 12
\begin{ex}
  Một hình vuông có chu vi là $\frac{8}{3}\text{ m}$. Diện tích hình vuông đó có mẫu số là bao nhiêu khi rút gọn tối giản?
  \shortans{9}
  \loigiai{
    Cạnh hình vuông là $\frac{8}{3} : 4 = \frac{2}{3}\text{ m}$. Diện tích là $\frac{2}{3} \times \frac{2}{3} = \frac{4}{9}\text{ m}^2$. Mẫu số là 9.
  }
\end{ex}

% Câu 13
\begin{ex}
  Một chai dầu chứa $\frac{3}{4}$ lít dầu. Hỏi $8$ chai dầu như thế chứa tất cả bao nhiêu lít dầu?
  \shortans{6}
  \loigiai{
    8 chai chứa số lít dầu là: $\frac{3}{4} \times 8 = 6$ (lít).
  }
\end{ex}

% Câu 14
\begin{ex}
  Có $12\text{ kg}$ đường đem chia đều vào các túi, mỗi túi chứa $\frac{3}{4}\text{ kg}$. Hỏi chia được tất cả bao nhiêu túi đường?
  \shortans{16}
  \loigiai{
    Số túi đường là: $12 : \frac{3}{4} = 12 \times \frac{4}{3} = 16$ (túi).
  }
\end{ex}

% Câu 15
\begin{ex}
  Một lớp học có $40$ học sinh. Trong đó có $\frac{1}{4}$ số học sinh giỏi và $\frac{1}{2}$ số học sinh khá. Hỏi lớp có bao nhiêu học sinh giỏi và khá?
  \shortans{30}
  \loigiai{
    Phân số chỉ học sinh giỏi và khá là $\frac{1}{4} + \frac{1}{2} = \frac{3}{4}$. Số học sinh là $40 \times \frac{3}{4} = 30$ học sinh.
  }
\end{ex}

% Câu 16
\begin{ex}
  Tìm giá trị của biểu thức: $(\frac{1}{2} + \frac{1}{3}) \times 12$.
  \shortans{10}
  \loigiai{
    $\frac{1}{2} \times 12 + \frac{1}{3} \times 12 = 6 + 4 = 10$.
  }
\end{ex}

% Câu 17
\begin{ex}
  Tìm giá trị của biểu thức: $(\frac{3}{4} - \frac{1}{6}) \times 24$.
  \shortans{14}
  \loigiai{
    $\frac{3}{4} \times 24 - \frac{1}{6} \times 24 = 18 - 4 = 14$.
  }
\end{ex}

% Câu 18
\begin{ex}
  Tính nhanh: $\frac{2024 \times 2025 - 1}{2023 \times 2025 + 2024}$. Kết quả là bao nhiêu?
  \shortans{1}
  \loigiai{
    Ta có $2024 \times 2025 - 1 = (2023 + 1) \times 2025 - 1 = 2023 \times 2025 + 2024$. Tử số bằng mẫu số nên kết quả bằng 1.
  }
\end{ex}

% Câu 19
\begin{ex}
  Một tấm vải dài $36\text{ m}$. Người ta cắt lấy $\frac{2}{3}$ tấm vải để may quần áo. Phần vải còn lại dài bao nhiêu mét?
  \shortans{12}
  \loigiai{
    Phần vải còn lại là $1 - \frac{2}{3} = \frac{1}{3}$ tấm vải, dài $36 \times \frac{1}{3} = 12\text{ m}$.
  }
\end{ex}

% Câu 20
\begin{ex}
  Tìm số nghịch đảo của phân số $\frac{5}{15}$. Kết quả là bao nhiêu?
  \shortans{3}
  \loigiai{
    Rút gọn $\frac{5}{15} = \frac{1}{3}$. Phân số nghịch đảo là $\frac{3}{1} = 3$.
  }
\end{ex}

% Câu 21
\begin{ex}
  Tính: $15 \times \frac{4}{5} : 3$. Kết quả là bao nhiêu?
  \shortans{4}
  \loigiai{
    $15 \times \frac{4}{5} = 12$. $12 : 3 = 4$.
  }
\end{ex}

% Câu 22
\begin{ex}
  Tính: $\frac{1}{2} : \frac{1}{4} : \frac{1}{8}$. Kết quả là bao nhiêu?
  \shortans{16}
  \loigiai{
    $\frac{1}{2} \times 4 \times 8 = 2 \times 8 = 16$.
  }
\end{ex}

% Câu 23
\begin{ex}
  Hai vòi nước cùng chảy vào một bể không có nước. Trong một giờ vòi một chảy được $\frac{1}{3}$ bể, vòi hai chảy được $\frac{1}{4}$ bể. Hỏi sau bao nhiêu giờ thì cả hai vòi chảy đầy $7$ bể như thế?
  \shortans{12}
  \loigiai{
    Trong 1 giờ hai vòi chảy được: $\frac{1}{3} + \frac{1}{4} = \frac{7}{12}$ bể. Để chảy đầy 7 bể cần: $7 : \frac{7}{12} = 12$ giờ.
  }
\end{ex}

% Câu 24
\begin{ex}
  Tính giá trị: $(1 - \frac{1}{2}) \times (1 - \frac{1}{3}) \times (1 - \frac{1}{4}) \times (1 - \frac{1}{5})$. Phân số tối giản có mẫu số là bao nhiêu?
  \shortans{5}
  \loigiai{
    Biểu thức bằng $\frac{1}{2} \times \frac{2}{3} \times \frac{3}{4} \times \frac{4}{5} = \frac{1}{5}$. Mẫu số là 5.
  }
\end{ex}

% Câu 25
\begin{ex}
  Tìm một số biết rằng $\frac{3}{5}$ của số đó bằng $45$. Số đó là bao nhiêu?
  \shortans{75}
  \loigiai{
    Số cần tìm là $45 : \frac{3}{5} = 45 \times \frac{5}{3} = 75$.
  }
\end{ex}
